\documentclass{article}
\usepackage{ISBI_template-master/spconf,amsmath,amssymb,graphicx}
\usepackage{booktabs,longtable,array,enumitem}
\usepackage[hidelinks]{hyperref}
\newcolumntype{P}[1]{>{\raggedright\arraybackslash}p{#1}}
\newcommand{\resultplot}[3]{%
  \begin{figure}[p]
  \centering
  \includegraphics[width=0.96\textwidth,height=0.75\textheight,keepaspectratio]{../plots/#1/#2.png}
  \caption{#3}
  \end{figure}%
}
\setlength{\parindent}{0pt}
\setlength{\parskip}{5pt}
\setlist[itemize]{nosep,leftmargin=1.4em}
\pagestyle{plain}
\title{Deep Double Descent Reproduction}
\name{Assignment 2: selected paper, replication track}
\address{Advanced Topics in Deep Learning, University of Copenhagen}
\begin{document}
\onecolumn
\maketitle



\begin{abstract}
In this paper, we try to reproduce the findings from 'Deep Double Descent: Where Bigger Models and More Data Hurt'. We find that we overall


\end{abstract}


\subsection{Methodology}












\section{Every figure and its experiment / compute dependency}
The table enumerates Figures 1--29 without treating 29 figures as 29 runs.
$W$ is an unresolved width count, $N$ a sample-count grid size, $D$ a
weight-decay count, and $s$ an unresolved repeat count. Values labeled R are
verified coordinates, not verified execution logs. Add overhead $e_i$ to
all training expressions. Where a figure reuses trajectories, its extra
\emph{training} cost is zero only if those trajectories and predictions were
saved; analysis/evaluation is still chargeable.
\begingroup\small
\setlength{\tabcolsep}{4pt}\renewcommand{\arraystretch}{1.13}
\begin{longtable}{@{}P{0.32in}P{3.47in}P{2.9in}@{}}
\toprule
Fig. & Experiment, varied quantities, fixed conditions and evidence & GPU-hours / reuse and remaining uncertainty\\
\midrule\endfirsthead
\toprule Fig. & Experiment and evidence & GPU-hours / reuse and uncertainty\\\midrule\endhead
1 & CIFAR-10, ResNet18; width sweep, 15\% noise, augmentation, Adam, 4k epochs. Train/test error and epoch-colored curves. Five trials (P, B.4); $k=1,\ldots,64$ (R). & $64\times5=320$ runs: $320A$. At 0.01/0.1 s: 1,390/13,902 GPU-h, before overhead. Figure 1 uses noisy-distribution test error.\\
2 & Train/test heatmaps versus width and epoch for the same 15\% noise ResNet family (P). & Reuse Figure 1 trajectories if identical: zero additional training. Heatmap cells are not separate training runs. Exact repeat aggregation unverified.\\
3 & IWSLT de--en, 4k versus 18k sentence pairs; Transformer width versus held-out perplexity (P). Released $d=8,16,\ldots,200$, 25 distinct widths each; 25+27 CSV rows (R). & $50T$ for distinct configurations, or $52T$ retaining all released rows as separate runs. At 0.1/0.5 s: 111--116 / 556--578 GPU-h. Two duplicate widths in the 18k file need provenance.\\
4 & ResNet18, CIFAR-10 and CIFAR-100; width and label-noise sweeps, augmentation, Adam, 4k epochs (P). Appendix A lists CIFAR-10 noise 0/5/10/15/20\%; CIFAR-100 0/10/20\%. & $64s(5+3)A=512sA$ if all listed conditions use the released 64-width grid (A). R verifies 64 widths in each base archive, not the entire condition cross product. Subtract exact overlaps with Figure 1; do not assume five seeds everywhere.\\
5 & CIFAR-10 CNN; widths, augmentation on/off, noise 0/10/20\%; SGD 500k steps (P). & $S\sum_{a,p}W_{a,p}s_{a,p}$. R: no-augmentation archive has 65 widths; augmented 10\% has 76 and 20\% has 60. Other grids/repeats need resolution. Figures 25/27 can reuse matching runs.\\
6 & CIFAR-10 CNN, clean labels, no augmentation; SGD 500k steps versus Adam 4k epochs across width (P). & R: 65 widths in each optimizer archive. Conditional one-repeat cost $65S+65A$; SGD clean runs may overlap Figure 5. Adam is separately trained.\\
7 & CIFAR-100 CNN, no label noise or augmentation; model-wise descent (P). Five trials in B.4; Figure 20 specifies one million steps. R: 64 widths. & $64\times5L=320L$ if all trials cover that grid: 889/8,889 GPU-h at 0.01/0.1 s. Do not cost this at 500k steps.\\
8 & Transformer width on full IWSLT de--en (160k pairs) and WMT en--fr (200k pairs), 80k updates; test perplexity (P). R: 64 and 62 rows/widths. & $64T_{\rm IWSLT}+62T_{\rm WMT}$. At common 0.1/0.5 s: 280/1,400 GPU-h. Different vocabulary/length distributions require separate timings.\\
9 & CIFAR-10 ResNet18 with 20\% noise and augmentation, Adam; three width slices and the width/epoch heatmap (P). & Reuse matching Figure 4 trajectories. Three displayed slices are not an additional three runs if already present.\\
10 & Epoch-wise curves: ResNet18 on CIFAR-10 and CIFAR-100; width-128 CNN on CIFAR-10, with noise levels (P). ResNet uses Adam; CNN uses SGD. & Reuse exact matching runs from Figures 4/5; otherwise one $A$ or $S$ per missing dataset/noise/width/seed tuple. $k=128$ is outside the base 1--64 ResNet archive; do not assume it is covered.\\
11(a) & CIFAR-10 CNN width versus sample size, augmentation, 10\%/20\% noise. Five trials for bottom panel, with initialization and subset randomness (P). & $S\sum_p W_pN_ps_p$, less exact reuse with Figure 12. Panel-specific grids/seeds not fully verified. B.4 says 500k ``epochs'', conflicting with the step-based protocol.\\
11(b) & IWSLT sample-count curves at $d=64$ and $80$ (P). R: each file has 59 counts, 4k through 120k in increments of 2k. & $118T$, or 262/1,311 GPU-h at 0.1/0.5 s. Four $(n,d)$ coordinates overlap Figure 3, but matching seeds/subsets are not established.\\
12 & CIFAR-10 CNN, 20\% noise; width/sample heatmap and three slices (P). R: 42 widths and nine sample sizes. & $42\times9=378$ configurations per seed, subject to valid cells: $378sS$. At 0.01/0.1 s: 525/5,250 GPU-h per seed. Reuse matching Figure 11(a) cells only once.\\
13 & Parameter count versus width/embedding dimension for CNN, ResNet and Transformer (P). & Zero training GPU-hours; instantiate/count on CPU. Transformer counts depend on vocabulary and embedding sharing, still unresolved.\\
14 & Fashion-MNIST random Fourier features: feature-count/sample-count heatmap and model-wise slice (P, Appendix D). & $\sum_{d,n,s}h_{\rm solve}(d,n,s)$. CPU-only implementation has zero GPU-h but nonzero CPU/RAM cost; GPU solve must be timed. Exact grid/solver/tolerance/seeds unverified.\\
15 & Sample-wise random-feature curve at 1,000 features (P). & Reuse the $d=1000$ slice of Figure 14 if present; otherwise additional fits for every sample count. Not a new deep-network sweep.\\
16 & CIFAR-10 ResNet18, 20\% noise/augmentation; Adam; constant, inverse-square-root and dynamic schedules. Legends show initial LR $10^{-4},10^{-3}$; axes extend to $10^6$ iterations (P, visual). & Six displayed settings per repeat. A full one-million-update horizon gives $6sL$; this horizon is a planning interpretation of axes, not a verified completion log. Width, dynamic-drop settings and actual terminal steps unresolved.\\
17 & Same task; SGD without momentum. Constant LR $10^{-4},10^{-3},10^{-2}$; inverse-square-root and dynamic LR $10^{-3},10^{-2},10^{-1}$ (P, visual). & Nine displayed settings: $9sL$ under the same full-horizon assumption. Reuse only configurations identical to other image runs; nominal SGD alone does not establish reuse.\\
18 & Same task; SGD with momentum 0.9. Each of three schedules shows LR $10^{-4},10^{-3},10^{-2}$ (P, visual). & Nine displayed settings: $9sL$ under the full-horizon assumption. Together Figures 16--18 have 24 settings, not merely nine optimizer/schedule pairs.\\
19 & Clean CIFAR-100 ResNet18, augmentation, Adam 4k epochs; width/epoch heatmaps and optimal early-stopping envelope (P). & Reuse clean CIFAR-100 runs from Figure 4. Early-stopping envelope can be computed from saved metrics; no retraining. This setting can retain model-wise double descent even with optimal early stopping.\\
20 & Clean CIFAR-100 CNN, no augmentation, SGD one million steps; dynamics/early stopping (P). & Same experiment as Figure 7: no additional training with saved histories.\\
21 & CIFAR-10 CNN, 10\% noise, augmentation, SGD 500k; width and weight decay $0,10^{-4},5\times10^{-4}$ (P, visual). & $S\sum_\lambda W_\lambda s_\lambda$; $180sS$ if all three use the 60-width decay grid (A). Zero-decay runs may already exist. R: decay directory listed only \texttt{ks}, no \texttt{Mlist}.\\
22 & CIFAR-10 ResNet18, 20\% noise/augmentation; regularization strength from 0 to 0.1 versus time, SGD inverse-square-root (P, E.2.2). & $\sum_{\lambda,s}U_{\lambda,s}t_{\lambda,s}/3600$. Exact decay grid, fixed width and horizons unresolved. A similarly named released archive says Adam, so it is not verified as this figure's data. Instability/reruns are an additional cost.\\
23 & Early-stopping/dynamics for IWSLT 4k/18k width sweeps (P). & Reuse Figure 3 training if epoch histories are saved. Released CSVs contain endpoint metrics only; they cannot reconstruct these trajectories.\\
24 & Early-stopping/dynamics for IWSLT 160k and WMT 200k width sweeps (P). & Reuse Figure 8 trajectories. Endpoint CSVs alone are insufficient; missing histories require new full runs or locating additional released logs.\\
25 & CIFAR-10 CNN, 10\% noise, augmentation, SGD inverse-square-root; width/epoch heatmaps (P, caption). & Reuse matching Figure 5/27 runs; conditional released 76-width family costs $76sS$ if generated from scratch. Appendix A's pointer for this figure conflicts with its caption.\\
26 & CIFAR-10 25k, clean labels; adversarial ResNet18 width sweep at $L_2$ radii 0.5 and 1.0, ten-step PGD. SGD LR 0.1 for 400 epochs then 0.01 for 400 (P). & $2WsP$ plus robust evaluation attacks. R: an \texttt{adv-mcnn} archive has 36 widths, but its architecture name conflicts with the caption. $72sP$ is conditional, not a verified ResNet run count.\\
27 & Wider augmented CIFAR-10 CNNs, noise conditions, SGD 500k; extension of Figure 5(b) (P). R: larger-model archives include 59/60 widths for 0/20\% noise. & Only widths/configurations absent from Figure 5 incur additional $S$ runs. Adding the full archive sizes would double-count shared small widths. Ten-percent family has its own 76-width list.\\
28 & CIFAR-10 ResNet18, 15\% noise; plurality-vote ensembles of five, with independently sampled corruption for each member; Adam 4k (P). & Reuse Figure 1 members only if independent-noise semantics and predictions match. Otherwise $5WA$ new training plus inference. Aggregate errors cannot reconstruct ensemble votes.\\
29 & CIFAR-10 CNN, clean labels, no augmentation; five-member ensemble, SGD (P, caption). & $5WsS$ if 500k is used, minus reusable independent members. Width grid/repeats not established. Caption points to Figure 7, which is CIFAR-100; do not reuse that different dataset.\\
\bottomrule
\end{longtable}
\endgroup

\section{Translation: verified grids and defensible cost arithmetic}
The six released tables contain the following exact coordinates. These
counts are verified from CSV contents, not inferred from axis tick marks.
They do not establish how many seeds originally produced each row.
\begin{center}\small
\begin{tabular}{@{}llll@{}}
\toprule
Released table / task & Grid & Rows & Unique within table\\\midrule
\texttt{nlp-model-dd.csv} / IWSLT & $d=8,16,\ldots,512$ & 64 & 64\\
\texttt{dd-translation-model-french.csv} / WMT & $d=8,16,\ldots,496$ & 62 & 62\\
\texttt{df4k.csv} / IWSLT 4k & $d=8,16,\ldots,200$ & 25 & 25\\
\texttt{df18k.csv} / IWSLT 18k & same grid & 27 & 25\\
\texttt{dd-translation-samples-aug26-small.csv} & $d=64$, 59 sample counts & 59 & 59\\
\texttt{dd-translation-samples-aug27-size80.csv} & $d=80$, 59 sample counts & 59 & 59\\\midrule
Total & & 296 & 294\\\bottomrule
\end{tabular}
\end{center}
The 18k table repeats widths 56 and 112. Without run IDs these could be
reruns, separate seeds or accidental duplication; they must not be silently
averaged or treated as independent trials. Four coordinates, $(n,d)$ with
$n\in\{4000,18000\}$ and $d\in\{64,80\}$, appear in both small-data width
sweeps and sample sweeps. If their complete configuration, seed and subset
match, 294 within-table unique settings reduce to 290 globally unique
configurations. Otherwise budget separately. Retaining all published rows
without reuse gives 296 runs, before any new repeats.

For one GPU per run and the stated 80,000-update protocol:
\[
 H_{\rm language}=\frac{80000}{3600}\sum_i t_i+E.
\]
At a uniform hypothetical $t_i=0.1$ seconds, the 290--296 accounting cases
cost \textbf{644.44--657.78 GPU-h}, or \textbf{\$2,249.11--\$2,295.64} at the
historical scenario rate. At $0.5$ seconds they cost
\textbf{3,222.22--3,288.89 GPU-h}, or \textbf{\$11,245.56--\$11,478.22}.
These include neither overhead nor repeated-seed replication. Three fresh
seeds multiply the chosen unique-configuration budget by three. They are
not claims that either timing is achievable on a B200.

Even the 50-configuration small-data Figure 3 sweep costs 111.11 GPU-h at
0.1 s/update and 555.56 GPU-h at 0.5 s/update. Thus 120 GPU-h could cover
that subset only if the average update took approximately 0.108 s or less,
with no overhead and only one repeat. It cannot be called the cost of all
translation experiments.

\textbf{Metric/provenance caveats discovered in the released data.}
The model-wise notebooks label the vertical axis cross-entropy, while the
paper captions call it perplexity. Some CSVs have identical train-error and
test-error columns, and the 18k error column uses a different numerical scale
from several other tables. The notebook legend also reverses the language
directions relative to the paper. These discrepancies do not alter row
counts, but the files must not be used uncritically as validated metrics.
Record loss units and log base before transforming to perplexity; use clean
held-out negative log-likelihood rather than the label-smoothed training
objective for a newly implemented evaluation.

\section{Image and auxiliary experiments: what can be priced now}
\textbf{The Figure 1 grid alone:} 64 widths and five trials imply 320 runs.
At 4,000 epochs and 391 batches/epoch that is 500,480,000 updates.
An assumed sweep-average 0.01 s/update costs 1,390.22 GPU-h and
\$4,851.88; at 0.1 s/update it costs 13,902.22 GPU-h and \$48,518.76.
These are conditional examples, not measured requirements, but they expose
why seven widths on a reduced dataset cannot price the original figure.

\textbf{The Figure 12 sample/width grid:} the nine released sample sizes are
3,125; 4,419; 6,250; 8,839; 12,500; 17,678; 25,000; 35,355; and 50,000.
Its 42 widths give 378 settings per repeat if every cell is required and
valid. With 500k updates each, $t=0.01$ s implies 525 GPU-h
(\$1,832.25); $t=0.1$ s implies 5,250 GPU-h (\$18,322.50).
The five-repeat statement for Figure 11(a)'s lower panel is not evidence
that the entire Figure 12 grid has five repeats.

\textbf{Optimizer/schedule checks:} the 24 visually identified settings in
Figures 16--18, each extended to one million updates, imply $24L$ per
repeat: 66.67/666.67 GPU-h at 0.01/0.1 s. Width and dynamic-drop parameters
still need recovery. Learning-rate curves that stop early in the plot do
not establish a common completed horizon.

\textbf{Random features and parameter counting:} CPU-only choices require
zero GPU-hours, but dense feature matrices and least-squares factorizations
can dominate RAM and CPU time. Exact dimensions, numerical method and
regularization/gradient-flow stopping must be specified and timed. Do not
reserve a GPU merely because the previous cost table used one as a proxy.
A large solve belongs on suitable remote compute.

\textbf{Adversarial training:} counting 800 epochs alone misses ten PGD
steps inside each update and separate robust evaluation. The ambiguous
36-width released archive would imply 72 runs for two radii only after its
architecture and condition mapping are resolved. It is not legitimate to
price this as two ordinary image sweeps.

\textbf{Ensembles:} five members per width, with the specified independent
corruptions for Figure 28. Preserve per-example predictions. The paper's
mean/standard-deviation curves and a CSV of aggregate losses cannot recover
plurality-vote ensemble errors. If members or predictions are missing,
training or inference is additional work.

\textbf{No grand total is added across figure rows.} Many rows are views of
shared trajectories, while several grids and seed counts remain unknown.
A defensible full total requires a union of run configurations, not the
sum of all displayed figure costs. The verified grids already show that the
full paper is materially larger than the previous proposed course subset.

\section{Conflicts that must remain visible}
\begin{enumerate}[leftmargin=1.5em,itemsep=3pt]
\item Appendix B.1 says ``80 gradient steps'' for Transformers, whereas
Section 4 and Figure 8 say 80k. This audit uses 80,000 and records the conflict.
\item Appendix B.4 says 500k \emph{epochs} for Figure 11(a); the main SGD
protocol is 500k \emph{updates}. Costs here use updates provisionally, not
500k full dataset passes. Original run configurations are needed to settle it.
\item Appendix A has figure mappings inconsistent with captions, including
Figure 25 and Figure 21's noise level. Caption-specific settings are used
provisionally: Figure 21 is 10\% noise; Figure 25 is augmented SGD at 10\%.
\item Figure 29's reference to Figure 7 conflicts on CIFAR-10 versus
CIFAR-100. They are not counted as the same training job.
\item The README calls translation optimization SGD; Appendix A says Adam
and B.3 references the Transformer recipe. This audit follows the paper
provisionally and does not claim released training code confirms it.
\item The adversarial archive says \texttt{mcnn}; Figure 26 says ResNet18.
The decay archive says Adam while E.2.2 specifies SGD. Directory names are
not enough to resolve either mapping.
\item Figure 1 evaluates against a noisy test distribution; other figures
use clean test targets. For ten classes and noise probability $p$, clean
error $e$ gives expected noisy error $p+(1-10p/9)e$. The inspected plotting
notebook instead uses an expression with the opposite sign on its final
$ep/9$ term. Check against original data before claiming exact curve recovery.
\item Optimal early stopping in published plots is an oracle envelope over
test trajectories. A new practical early-stopping experiment must use a
validation split. Figure 19 is an explicit exception to the broad appendix
heading claiming early stopping does not exhibit double descent.
\end{enumerate}

\section{A costable course subset, without concealing its reductions}
A useful cross-modality candidate is one image width/dynamics experiment,
one translation width/sample experiment and one independent ablation. This
is a proposed menu, not an approved compute launch or a claim that it
reproduces every figure.

\textbf{Image contribution:} seven pilot-selected ResNet widths, CIFAR-10
50k, 15\% noise, augmentation, full 4k epochs, three seeds: 21 runs.
It targets Figures 1/2 but uses fewer widths and repeats. A clean-label
ablation at three of those widths adds nine runs. Both epoch-wise and
model-wise plots use the same saved histories. Total: $30A$, or
130.33/1,303.33 GPU-h at hypothetical 0.01/0.1 s per update.

\textbf{Translation contribution:} six pilot-selected dimensions (multiples
of eight) at both 4k and 18k pairs, three seeds: 36 runs at 80k updates.
It targets Figures 3/23 and a sparse cross-section of Figure 11(b), with
explicitly reduced grid density. Total: $36T$, or 80/400 GPU-h at
hypothetical 0.1/0.5 s per update. A separate WMT comparison is not included;
adding six WMT widths and three seeds costs another $18T_{\rm WMT}$.

The image sweep, translation sweep and clean-label ablation give three
concrete experimental contributions without pretending image-only work
covers the text results. At the faster illustrative timings, the combined
$30A+36T$ is 210.33 GPU-h or \$734.06 before overhead; doubling that amount
is \$1,468.13. At the slower illustrative timings it is 1,703.33 GPU-h or
\$5,944.63 before overhead. \textbf{Even this subset cannot yet be promised
under \$1,000.} Fewer seeds, fewer widths, smaller image subsets or shorter
horizons require a new explicit scope decision; stopping before interpolation
can remove the phenomenon being tested. Pilot results must determine the
choice rather than retroactively justifying an attractive price.


\section{Completed reduced vision results}
After the initial audit, our group completed 229 single-seed vision cells.
The Adam runs end at epoch 400 and the SGD CNN runs end at 50,000 optimizer
steps. The authors generally used denser width grids, 4,000 Adam epochs,
500,000 CIFAR-10 SGD updates, or one million clean CIFAR-100 SGD updates.
The figures below therefore document partial reproductions of the numbered
paper figures. They do not establish a full replication of any numbered
figure. Our training time is recorded as epochs or optimizer steps; the
authors' released arrays often supply only a measurement index. For CNN
panels near 50,000 steps, the authors' step scale uses an unverified
256-step-per-measurement estimate. Our single-seed results have no inferred
error bars. The data and remaining gaps are detailed in
\texttt{../EXPERIMENTS.md}.

\subsection{Our measurements beside the released measurements}
Each image pairs our results with a plot reconstructed from the authors'
released numeric data. The image labels identify the source of each panel;
matching visual layout does not imply equal training horizons.

\resultplot{comparison}{figure_01}{Figure 1 adaptation: CIFAR-10 ResNet18 at 10\% label noise. Our ten widths and 400 epochs cover one seed; the paper uses 15\% noise, more widths, five trials, and 4,000 epochs.}
\resultplot{comparison}{figure_02}{Figure 2 adaptation: train and test error heatmaps at 10\% noise. The published measurement positions and our recorded epochs are displayed separately.}
\resultplot{comparison}{figure_04}{Figure 4 adaptation: CIFAR-10 and CIFAR-100 ResNet18 noise curves. Our 0/10/20\% noise cells end at epoch 400; the published CIFAR-10 set also includes 5/15\% noise.}
\resultplot{comparison}{figure_05}{Figure 5 subset: augmented and unaugmented CIFAR-10 CNN curves. Our runs end at 50,000 SGD steps; the authors' step mapping near that horizon is estimated.}
\resultplot{comparison}{figure_06}{Figure 6 subset: clean, unaugmented CIFAR-10 CNN with SGD and Adam. The published endpoints are from full recorded histories; ours stop at 50,000 steps or 400 epochs.}
\resultplot{comparison}{figure_07}{Figure 7 subset: clean CIFAR-100 CNN. The published five-trial mean and spread are shown beside our one-seed 50,000-step result at eleven widths.}
\resultplot{comparison}{figure_09}{Figure 9 adaptation: 20\%-noise CIFAR-10 ResNet18 trajectories and width-by-time heatmap. Our run ends at epoch 400.}
\resultplot{comparison}{figure_10}{Figure 10 subset: width-128 CIFAR-10 CNN histories at three noise levels. Our curves use recorded steps; the published step coordinate is inferred.}
\resultplot{comparison}{figure_11a}{Figure 11(a) subset: CIFAR-10 CNN sample-size curves. The 20\%-noise comparison uses three sample sizes and eleven widths; our 10\%-noise results appear separately because matching published subset data were unavailable.}
\resultplot{comparison}{figure_12}{Figure 12 subset: 20\%-noise CIFAR-10 CNN heatmaps and sample-size slices at three sample sizes and eleven widths, with a common color scale.}

\clearpage
\subsection{Our measurements alone}
The following panels use the same plot definitions and our same measured
cells as the paired images above. They make our observed curves readable
without the additional published panels. The experimental limits stated
above apply to every image in this subsection.

\resultplot{ours}{figure_01}{Our Figure 1 adaptation: 10\%-noise CIFAR-10 ResNet18 endpoints and training-time curves through epoch 400.}
\resultplot{ours}{figure_02}{Our Figure 2 adaptation: 10\%-noise CIFAR-10 ResNet18 test and train heatmaps through epoch 400.}
\resultplot{ours}{figure_04}{Our Figure 4 adaptation: CIFAR-10 and CIFAR-100 ResNet18 curves at 0/10/20\% noise.}
\resultplot{ours}{figure_05}{Our Figure 5 subset: augmented and unaugmented CIFAR-10 CNN curves through 50,000 SGD steps.}
\resultplot{ours}{figure_06}{Our Figure 6 subset: clean, unaugmented CIFAR-10 CNN endpoints for SGD and Adam.}
\resultplot{ours}{figure_07}{Our Figure 7 subset: one-seed clean CIFAR-100 CNN endpoints after 50,000 SGD steps.}
\resultplot{ours}{figure_09}{Our Figure 9 adaptation: 20\%-noise CIFAR-10 ResNet18 trajectories and width-by-epoch heatmap.}
\resultplot{ours}{figure_10}{Our Figure 10 subset: width-128 CIFAR-10 CNN test and train histories through 50,000 steps.}
\resultplot{ours}{figure_11a}{Our Figure 11(a) subset: three CIFAR-10 CNN sample sizes at 10\% and 20\% noise.}
\resultplot{ours}{figure_12}{Our Figure 12 subset: 20\%-noise CIFAR-10 CNN sample-size heatmap and slices, using the same color scale as the paired comparison.}
\end{document}
